\section{Presheaves}
\label{section-presheaves}

\noindent
Let $\mathcal{C}$ be a category.
A {\it presheaf of sets} is a contravariant functor $\mathcal{F}$
from $\mathcal{C}$ to $\textit{Sets}$ (see Categories, Remark
\ref{categories-remark-functor-into-sets}).
So for every object $U$ of $\mathcal{C}$ we have a set
$\mathcal{F}(U)$. The elements of this set are called
the {\it sections} of $\mathcal{F}$ over $U$. For every morphism
$f : V \to U$ the map $\mathcal{F}(f) : \mathcal{F}(U) \to \mathcal{F}(V)$
is called the {\it restriction map} and is often denoted
$f^\ast : \mathcal{F}(U) \to \mathcal{F}(V)$. Another way
of expressing this is to say that $f^*(s)$ is the {\it pullback}
of $s$ via $f$. Functoriality means that $g^* f^* (s) = (f \circ g)^*(s)$.
Sometimes we use the notation $s|_V := f^\ast(s)$.
This notation is consistent with the notion of restriction
of functions from topology because if $W \to V \to U$
are morphisms in $\mathcal{C}$ and $s$ is a section of
$\mathcal{F}$ over $U$ then $s|_W = (s|_V)|_W$ by the
functorial nature of $\mathcal{F}$. Of course we have to be
careful since it may very well happen
that there is more than one morphism $V \to U$ and it is
certainly not going to be the case that the corresponding
pullback maps are equal.

\begin{definition}
\label{definition-presheaves-sets}
A {\it presheaf of sets} on $\mathcal{C}$ is a contravariant
functor from $\mathcal{C}$ to $\textit{Sets}$. {\it Morphisms
of presheaves} are transformations of functors. The category
of presheaves of sets is denoted $\textit{PSh}(\mathcal{C})$.
\end{definition}

\noindent
Note that for any object $U$ of $\mathcal{C}$ the functor of
points $h_U$, see Categories, Example \ref{categories-example-hom-functor}
is a presheaf. These are called the {\it representable presheaves}.
These presheaves have the pleasing property that for any
presheaf $\mathcal{F}$ we have
\begin{equation}
\label{equation-map-representable-into-presheaf}
\Mor_{\textit{PSh}(\mathcal{C})}(h_U, \mathcal{F})
=
\mathcal{F}(U).
\end{equation}
This is the Yoneda lemma (Categories, Lemma \ref{categories-lemma-yoneda}).

\medskip\noindent
Similarly, we can define the notion of a presheaf of abelian groups,
rings, etc. More generally we may define a presheaf with values in a
category.

\begin{definition}
\label{definition-presheaf}
Let $\mathcal{C}$, $\mathcal{A}$ be categories.
A {\it presheaf} $\mathcal{F}$ on $\mathcal{C}$
with values in $\mathcal{A}$ is a contravariant
functor from $\mathcal{C}$ to $\mathcal{A}$,
i.e., $\mathcal{F} : \mathcal{C}^{opp} \to \mathcal{A}$.
A {\it morphism} of presheaves $\mathcal{F} \to \mathcal{G}$
on $\mathcal{C}$ with values in $\mathcal{A}$ is a transformation
of functors from $\mathcal{F}$ to $\mathcal{G}$.
\end{definition}

\noindent
These form the objects and morphisms of the category of presheaves
on $\mathcal{C}$ with values in $\mathcal{A}$.

\begin{remark}
\label{remark-big-presheaves}
As already pointed out we may consider the category of
presheaves with values in any of the ``big'' categories
listed in Categories, Remark \ref{categories-remark-big-categories}.
These will be ``big'' categories as well and they will be
listed in the above mentioned remark as we go along.
\end{remark}
















